% ================================================================
% SECTION 14: CONFIGURATION SYSTEM
% ================================================================
\section{Configuration System}
\label{sec:config}

\begin{center}
\code{bionetflux.utils.config\_manager} \quad(550 lines) \\
\code{bionetflux.problems.ks\_config\_manager} \quad(149 lines) \\
\code{bionetflux.problems.ooc\_config\_manager} \quad(168 lines)
\end{center}

All problem parameters, discretization settings, initial conditions,
forcing functions, and boundary-condition overrides are specified in a
TOML file.  The configuration system provides loading, merging with
defaults, function resolution, and validation.

% ------------------------------------------------------------------
\subsection{TOML File Structure}

A configuration file is organized into the following top-level sections:

\begin{lstlisting}[caption={Typical TOML layout (Keller-Segel example)}]
[problem]
name = "KS_Traveling_Wave"
neq = 2
problem_type = "ks"

[time_parameters]
T = 0.5
dt = 0.05

[physical_parameters.diffusion]
mu = 2.0
nu = 1.0

[physical_parameters.reaction]
a = 0.0
b = 1.0

[physical_parameters.chemotaxis]
chi = "constant"
dchi = "zeros"

[discretization]
n_elements = 10
tau = [1.0, 1.0]

[initial_conditions]
u = "zeros"
phi = "zeros"

[force_functions]
u = "zeros"
phi = "zeros"

[boundary_conditions]
B1_u = { type = "dirichlet", data = "zeros" }
\end{lstlisting}

% ------------------------------------------------------------------
\subsection{load\_toml\_config}

\begin{lstlisting}[language=Python, caption={TOML loader}]
def load_toml_config(config_file: str) -> Dict[str, Any]:
\end{lstlisting}

Tries \code{tomllib} (Python $\geq$3.11) first, then falls back to the
\code{tomli} package.

% ------------------------------------------------------------------
\subsection{FunctionResolver}

\begin{lstlisting}[language=Python, caption={FunctionResolver essentials}]
class FunctionResolver:
    def __init__(self): ...
    def resolve_function(self, func_spec) -> Callable: ...
    def resolve_boundary_function(self, func_spec, position) -> Callable: ...
    def parse_expression(self, expr_str, variables=['s','t']) -> Callable: ...
    def register_function(self, name, func): ...
    def register_function_library(self, library): ...
\end{lstlisting}

Function resolution follows a three-step cascade:
\begin{enumerate}
  \item If \code{func\_spec} is already callable, return it.
  \item If the string appears in the built-in registry, return the registered callable.
  \item Otherwise parse it as a SymPy expression via \code{sympy.sympify} $\to$ \code{sympy.lambdify}.
\end{enumerate}

\noindent Built-in common library (selected entries):

\begin{longtable}{p{4.5cm}p{9cm}}
\toprule
\textbf{Name} & \textbf{Definition ($s =$ space, $t =$ time)} \\
\midrule
\code{zeros} & $f(s,t) = 0$ \\
\code{ones} / \code{constant} & $f(s,t) = 1$ \\
\code{sin\_2pi} / \code{cos\_2pi} & $\sin/\cos(2\pi s)$ \\
\code{gaussian} & $\exp(-s^2)$ \\
\code{quadratic} & $s(1-s)$ \\
\code{step} & $H(s - 0.5)$ \\
\code{traveling\_wave\_u} & KS traveling-wave $u(s,t)$ \\
\code{traveling\_wave\_phi} & KS traveling-wave $\varphi(s,t)$ \\
\bottomrule
\end{longtable}

\code{resolve\_boundary\_function(func\_spec, position)} wraps the resolved
$f(s,t)$ into $g(t) = f(\text{position},\, t)$, which is the signature expected by
\code{Constraint.data\_function}.

% ------------------------------------------------------------------
\subsection{ParameterValidator}

\begin{lstlisting}[language=Python, caption={ParameterValidator}]
class ParameterValidator:
    def add_rule(self, param_path: str, rule: dict): ...
    def validate_config(self, config) -> List[str]: ...
\end{lstlisting}

Rules are keyed by dot-separated paths (e.g.\
\code{"physical\_parameters.diffusion.mu"}) and checked for type,
minimum/maximum, positivity, and required-ness.

% ------------------------------------------------------------------
\subsection{BaseConfigManager (ABC)}

\begin{lstlisting}[language=Python, caption={BaseConfigManager skeleton}]
class BaseConfigManager(ABC):
    def __init__(self, problem_type: str):
        self.function_resolver = FunctionResolver()
        self.validator = ParameterValidator()

    @abstractmethod
    def _setup_defaults(self): ...
    @abstractmethod
    def _setup_function_library(self): ...
    @abstractmethod
    def _setup_validation_rules(self): ...

    def load_config(self, config_file=None) -> Dict[str, Any]:
        """Load, merge with defaults, validate, resolve functions."""
\end{lstlisting}

The \code{load\_config} workflow:
\begin{enumerate}
  \item If a file is provided, load it with \code{load\_toml\_config}.
  \item Recursively merge with \code{self.defaults}.
  \item Run \code{ParameterValidator.validate\_config}; raise on errors.
  \item Call \code{\_resolve\_functions} to convert string specs in
        \code{initial\_conditions} and \code{force\_functions} to callables.
  \item Return the resolved configuration dictionary.
\end{enumerate}

% ------------------------------------------------------------------
\subsection{KSConfigManager and OoCConfigManager}

\begin{longtable}{p{4.5cm}p{4cm}p{5.5cm}}
\toprule
\textbf{Class} & \textbf{Problem type key} & \textbf{Additional function library} \\
\midrule
\code{KSConfigManager} & \code{"ks"} & \code{ks\_gaussian\_blob}, \code{ks\_step\_function}, \code{chemotactic\_gradient}, \code{attraction\_source} \\
\code{OoCConfigManager} & \code{"ooc"} & \code{ooc\_linear\_gradient}, \code{ooc\_exponential}, \code{nutrient\_gaussian} \\
\bottomrule
\end{longtable}

Both inherit from \code{BaseConfigManager}, define their own defaults (physical parameters matching the respective equation system), validation rules, and problem-specific functions.
